\documentclass[11pt,oneside]{book}

%------------------------------------------------
% Fonts
%------------------------------------------------

\usepackage{fontspec}
\setmainfont{TeX Gyre Pagella}

\usepackage{unicode-math}
\setmathfont{Latin Modern Math}

%------------------------------------------------
% Layout
%------------------------------------------------

\usepackage[margin=1.1in]{geometry}

\usepackage{microtype}

%------------------------------------------------
% Mathematics
%------------------------------------------------

\usepackage{amsmath}
\usepackage{amssymb}
\usepackage{amsthm}
\usepackage{mathtools}
\usepackage{mathrsfs}
\usepackage{bm}

%------------------------------------------------
% Figures
%------------------------------------------------

\usepackage{graphicx}
\usepackage{xcolor}
\usepackage{booktabs}
\usepackage{longtable}
\usepackage{array}
\usepackage{multirow}

%------------------------------------------------
% Hyperlinks
%------------------------------------------------

\usepackage[
colorlinks=true,
linkcolor=blue!50!black,
citecolor=blue!50!black,
urlcolor=blue!50!black
]{hyperref}

\usepackage[nameinlink]{cleveref}

%------------------------------------------------
% Index
%------------------------------------------------

\usepackage{imakeidx}
\makeindex

%------------------------------------------------
% Theorems
%------------------------------------------------

\numberwithin{equation}{chapter}

\theoremstyle{definition}
\newtheorem{definition}{Definition}[chapter]
\newtheorem{example}[definition]{Example}
\newtheorem{remark}[definition]{Remark}

\theoremstyle{plain}
\newtheorem{theorem}[definition]{Theorem}
\newtheorem{lemma}[definition]{Lemma}
\newtheorem{proposition}[definition]{Proposition}
\newtheorem{corollary}[definition]{Corollary}

%------------------------------------------------
% Operators
%------------------------------------------------

\DeclareMathOperator{\Hom}{Hom}
\DeclareMathOperator{\Obj}{Obj}
\DeclareMathOperator{\Aut}{Aut}
\DeclareMathOperator{\End}{End}
\DeclareMathOperator{\id}{id}
\DeclareMathOperator{\grad}{grad}

%------------------------------------------------
% Commands
%------------------------------------------------

\newcommand{\State}{\mathcal{S}}
\newcommand{\Theory}{\mathfrak{T}}
\newcommand{\Adm}{\mathsf{Adm}}
\newcommand{\Repair}{\mathcal{R}}
\newcommand{\Constraint}{\mathcal{C}}

%------------------------------------------------
% Title
%------------------------------------------------

\title{
The Evolution of RSVP\\[1em]
\Large
From Language to Plenum\\
The Development of a Constraint-Based
Research Program
}

\author{
Flyxion\\
Independent Researcher
}

\date{2026}

\begin{document}

\frontmatter

\maketitle

\tableofcontents

\chapter*{A Note on This Draft}
\addcontentsline{toc}{chapter}{A Note on This Draft}

This is a working compilation of the monograph in its current state of
development. Part I is drafted in full. Parts II through V, together with
the appendices, are presented here as structural outlines: each chapter is
named and its intended content described in prose, but the mathematical and
expository development itself remains to be written. The outline sections
are marked clearly at the start of each part so that drafted and undrafted
material are never confused with one another.

\mainmatter

\part{Foundations of Organization}

\chapter{Introduction}

\section{Why Another Theory?}

Every scientific theory begins by deciding what constitutes a fundamental
object. Classical mechanics begins with particles. General relativity begins
with spacetime geometry. Quantum mechanics begins with state vectors and
operators. Computer science often begins with symbols or computation.
Linguistics frequently begins with words, grammar, or meaning.

This book begins somewhere else.

It begins with the observation that every successful scientific description
must preserve distinctions.

Regardless of discipline, every theory partitions experience into states
which may be distinguished from one another. A physical experiment
distinguishes outcomes. A proof distinguishes valid from invalid reasoning.
A biological organism distinguishes self from environment. A compiler
distinguishes correct programs from incorrect ones. Language distinguishes
meanings through relations among symbols.

The particular objects differ. The underlying requirement does not.

The central question therefore becomes not

\begin{quote}
"What objects exist?"
\end{quote}

but instead

\begin{quote}
"How are distinctions preserved, transformed, repaired, and propagated
through time?"
\end{quote}

This seemingly modest change of perspective gradually alters the structure of
the entire research program.

Rather than treating particles, symbols, organisms, or economic agents as
primitive entities, they become stable patterns maintained by deeper
organizational constraints.

Objects become persistent trajectories.

Identity becomes maintained continuation.

Structure becomes preserved admissibility.

The remainder of this monograph documents how this perspective emerged over
several years of research.

\section{The Repository as Research Record}

Modern scientific research increasingly develops in public repositories.

Unlike traditional publication, where intermediate ideas disappear during
peer review, open repositories preserve the entire developmental history of
a research program. Incorrect approaches, abandoned terminology, failed
proofs, rewritten models, experimental software, and speculative
mathematics remain visible.

This historical record is unusually valuable.

The purpose of this volume is therefore not simply to summarize completed
results, but to reconstruct the conceptual evolution that produced them.

Many ideas presented later in their mature mathematical form first appeared
only as informal observations, software architectures, visualization
experiments, or attempts to solve practical engineering problems.

The reader should therefore expect recurring concepts to evolve rather than
appear fully formed.

This mirrors the actual process of scientific discovery.

The account given in this chapter is deliberately the simplest available
on-ramp into that history, organized around the two threads, language and
biology, that this book develops first. It is not the complete
conceptual history of the research program. That fuller history, in which
the physical theory introduced below turns out to be one of several
parallel developments rather than the trunk from which everything else
grew, is given in Appendix H and referenced again in Chapter 37.

\section{From Language to Physics}

The path traced in this section and the next is the simplest route into
the material, not the only one the research program actually took.
Presented here as a single line running from language to physics, it is
in fact one strand of a wider, branching development described in full in
Appendix H. That appendix will show physics, computation, and the theory
of observation developing as parallel realizations of a shared
ontological commitment, rather than one leading tidily to the next.
Nothing in the mathematics that follows depends on the simplified version;
it is offered because a single thread is easier to follow on a first
reading than a branching one.

The earliest investigations that eventually led to RSVP were not studies of
physics.

They were studies of language.

Questions concerning semantics, representation, reference, compression,
categorization, and symbolic interpretation repeatedly encountered the same
difficulty: meaning could not be identified with isolated symbols.

Meaning appeared instead to arise from structured relationships maintained
across changing contexts.

This observation gradually generalized.

The same mathematical ideas appeared within machine learning, compiler
design, distributed systems, physiological regulation, economic dynamics,
thermodynamics, and eventually theoretical physics.

Each discipline seemed to describe different phenomena while repeatedly
constructing remarkably similar mathematical machinery.

The research therefore shifted from studying individual disciplines toward
studying the mathematical structures common to all of them.

\section{A Shift in Perspective}

The central transition documented throughout this book may be summarized as
a sequence of conceptual replacements.

\begin{center}
\begin{tabular}{ll}
Objects & $\longrightarrow$ Processes \\
States & $\longrightarrow$ Trajectories \\
Representations & $\longrightarrow$ Constraints \\
Optimization & $\longrightarrow$ Repair \\
Static Structure & $\longrightarrow$ Continuation \\
Information & $\longrightarrow$ Admissibility
\end{tabular}
\end{center}

Each replacement enlarges the scope of the framework.

Rather than discarding earlier theories, the objective is to identify the
conditions under which those theories remain admissible approximations of a
more general mathematical description.

In this sense the framework is cumulative rather than revolutionary.

\section{The RSVP Program}

Relativistic Scalar Vector Plenum (RSVP) represents the physical expression
of this broader research program.

It proposes that the primitive description of physical reality is not a
collection of independent particles embedded within spacetime, but a
relativistic plenum whose scalar and vector fields evolve under local
constraints.

Observable objects correspond to dynamically maintained stable structures
within this field rather than irreducible building blocks.

Although this monograph discusses the physical implications of RSVP, its
primary objective is broader.

The physical theory is presented as one manifestation of a general
constraint-based mathematical language whose applications extend well beyond
physics.

It should be stated plainly here, rather than left for the closing
chapters: RSVP is not the foundation from which the rest of this book's
subject matter was derived. It is one branch of a broader research
program, organized around a shared reachability ontology, that also
produced independent computational and epistemic developments referenced
throughout Part IV and described fully in Appendix H. RSVP is given
first, and at greatest length, because it is the most mathematically
developed of these branches and because physics offers the cleanest tests
of the framework's central claims. That expository choice should not be
read as a claim of conceptual priority.

\section{How to Read This Book}

This book may be approached from several different backgrounds.

Readers primarily interested in theoretical physics may proceed directly to
the later chapters introducing the scalar-vector plenum.

Readers interested in artificial intelligence or computation will find that
many of the mathematical ideas originate within discussions of semantics,
software architecture, and admissible computation.

Readers interested in philosophy may view the work as an investigation into
the nature of identity, persistence, and structural realism.

Despite these differing motivations, all paths eventually converge upon the
same central idea:

stable reality is not defined by the existence of isolated objects, but by
the persistence of admissible relational structure.

Readers curious from the outset about how the physical, computational, and
epistemic material relate to one another, rather than encountering that
question only in Chapter 37 and the Epilogue, may wish to read Appendix H
early rather than last; it does not depend mathematically on anything
between here and there.

The remaining chapters develop this claim step by step, beginning where the
research itself began, with language and biology together.

\chapter{Two Lines of Evidence: Language and Biology}

